\section{Provenance and dominance reductions for the recorded source-pool frontiers}\label{app:frontiers}

Each modulus below is a separate existence guarantee.  Products are formed
only from disjoint input blocks.  After the direct moduli and the available
split products have been collected, a modulus $m$ is deleted when another
proved modulus $M$ in the same class satisfies $m\mid M$.  This is a genuine
dominance deletion: the expression producing an $M$-multiple is itself an
$m$-multiple.  It is not an operation that combines two different
expressions.  For instance, $21$ is deleted once $63$ is proved because
$21\mid63$, whereas separate guarantees $6$ and $8$ cannot be merged into a
$24$-guarantee without a separate proof.

The source pools used here consist of the direct theorems selected for this
summary together with products of the earlier displayed source-pool frontiers
over disjoint splits.  The lists below give the retained values and the
corresponding dominance deletions within those pools.  They are not intended
to be an exhaustive closure under every parity case split available elsewhere
in the manuscript, and no optimality claim is made.

\subsection{One and two inputs}

For one input,
\[
 F_1=O_1=\{1\},\qquad E_1=\{2\}.
\]
For two inputs, the universal guarantees are $2$ and $3$.  In the all-odd
class the additional modulus $4$ dominates $2$; in the all-even class the
moduli $6$ and $8$ dominate $2$ and $3$ as indicated below:
\[
 F_2=\{2,3\},\qquad O_2=\{3,4\},\qquad E_2=\{6,8\}.
\]
\[
\begin{array}{c|c}
\text{deleted modulus} & \text{retained dominator}\\ \hline
O_2:\ 2 & 4\\
E_2:\ 2 & 6\\
E_2:\ 3 & 6
\end{array}
\]

\subsection{Three inputs}

The universal direct constructions give $7,8,9,10,12$; the smaller split
products $2,3$ are dominated.  The parity-restricted direct improvements are
$14,18$ for odd inputs and $14,18,48$ for even inputs.  Thus
\[
\begin{aligned}
F_3&=\{7,8,9,10,12\},\\
O_3&=\{8,10,12,14,18\},\\
E_3&=\{10,14,18,48\}.
\end{aligned}
\]
The deletions are
\[
\begin{array}{c|c@{\qquad}c|c@{\qquad}c|c}
\multicolumn{2}{c}{F_3} & \multicolumn{2}{c}{O_3} & \multicolumn{2}{c}{E_3}\\
2&8 & 3&18 & 7&14\\
3&9 & 4&8  & 8&48\\
 &&7&14 & 9&18\\
 &&9&18 & 12&48\\
 &&     & 16&48
\end{array}
\]
where each left entry divides the adjacent retained entry.

\subsection{Four inputs}

The universal source pool contains the one-plus-three and two-plus-two split
products, together with the direct moduli $10,14,16,24,36$.  After dominance
reduction,
\[
 F_4=\{10,14,16,24,36\}.
\]
For four odd inputs the direct moduli $48$ and $60$ supplement the inherited
and split products, giving
\[
 O_4=\{14,36,48,60\}.
\]
For four even inputs the retained products and direct theorem are
\[
20=2\cdot10,\quad28=2\cdot14,\quad64=8\cdot8,
\quad72,\quad96=2\cdot48,
\]
so
\[
 E_4=\{20,28,64,72,96\}.
\]
The complete deletions from these source pools are
\[
\begin{array}{c|c@{\quad}c|c@{\quad}c|c}
\multicolumn{2}{c}{F_4} & \multicolumn{2}{c}{O_4} & \multicolumn{2}{c}{E_4}\\
4&16 & 8&48 & 36&72\\
6&24 & 9&36 & 48&96\\
7&14 & 10&60 &&\\
8&16 & 12&48 &&\\
9&36 & 16&48 &&\\
12&24 & 18&36 &&
\end{array}
\]

\subsection{Five inputs}

A disjoint $3+2$ split gives the universal candidates
\[
14,16,18,20,21,24,27,30,36,
\]
and the direct factorial theorem adds $120$.  Since
\[
18\mid36,\qquad20\mid120,\qquad24\mid120,\qquad30\mid120,
\]
the universal frontier is
\[
F_5=\{14,16,21,27,36,120\}.
\]

For five odd inputs, the $3\OO+2\OO$ products give
\[
24,30,32,36,40,42,48,54,56,72,
\]
and the direct theorems add $180,240$.  The retained values are
\[
O_5=\{32,42,54,56,72,180,240\},
\]
with deletions
\[
24\mid240,\quad30\mid180,\quad36\mid180,
\quad40\mid240,\quad48\mid240.
\]

For five even inputs, the $3\EE+2\EE$ products give
\[
60,80,84,108,112,144,288,384,
\]
and the direct theorem adds $480$.  Therefore
\[
E_5=\{84,108,112,288,384,480\},
\]
with
\[
60\mid480,\qquad80\mid480,\qquad144\mid288.
\]

\subsection{Six inputs}

For arbitrary inputs, the retained $3+3$ products are
\[
\begin{gathered}
49=7\cdot7,\quad56=7\cdot8,\quad63=7\cdot9,
\quad64=8\cdot8,\quad70=7\cdot10,\\
81=9\cdot9,\quad84=7\cdot12,\quad96=8\cdot12,
\quad100=10\cdot10,\quad108=9\cdot12.
\end{gathered}
\]
The direct six-input theorem adds $720$, giving
\[
F_6=\{49,56,63,64,70,81,84,96,100,108,720\}.
\]
All other candidates produced by the $1+5$, $2+4$, and $3+3$ splits are
removed by the following explicit divisibilities:
\[
\begin{array}{c|c@{\qquad}c|c@{\qquad}c|c@{\qquad}c|c}
14&56 &16&64 &20&100 &21&63\\
27&81 &28&56 &30&720 &32&64\\
36&108&42&84 &48&96  &72&720\\
80&720&90&720&120&720&144&720
\end{array}
\]

For six odd inputs, the retained products are
\[
\begin{gathered}
100=10\cdot10,\quad112=8\cdot14,\quad140=10\cdot14,
\quad168=12\cdot14,\quad192=4\cdot48,\\
196=14\cdot14,\quad216=12\cdot18,\quad252=14\cdot18,
\quad324=18\cdot18,
\end{gathered}
\]
together with $720$.  Thus
\[
O_6=\{100,112,140,168,192,196,216,252,324,720\}.
\]
The deleted candidates and one retained dominator for each are
\[
\begin{array}{c|c@{\qquad}c|c@{\qquad}c|c}
32&192 &42&168 &54&216\\
56&112 &64&192 &72&216\\
80&720 &96&192 &108&216\\
120&720&144&720&180&720\\
240&720&&&
\end{array}
\]

For six even inputs, the retained products are
\[
\begin{gathered}
100=10\cdot10,\quad140=10\cdot14,\quad196=14\cdot14,
\quad252=14\cdot18,\quad324=18\cdot18,\\
512=8\cdot64,\quad672=14\cdot48,\quad864=18\cdot48,
\quad960=2\cdot480,\quad2304=48\cdot48,
\end{gathered}
\]
together with $720$.  Hence
\[
E_6=\{100,140,196,252,324,512,672,720,864,960,2304\}.
\]
The complete deletion list is
\[
\begin{array}{c|c@{\qquad}c|c@{\qquad}c|c}
120&720 &160&960 &168&672\\
180&720 &216&864 &224&672\\
384&2304&432&864&480&960\\
576&2304&768&2304&&
\end{array}
\]

These source and deletion lists justify the displayed source-pool frontiers.
They do not exclude further nondominated moduli obtainable from other proved
case splits or from later strengthening arguments.
